\chapter*{不定积分每日一题}

\refer{《积分的方法与技巧》金玉明\ 顾新身\ 毛瑞庭\ /\ 编著}

\exercise{求积分$\int{x\arctan{x}\d x}$}

\teacher{刘嘉玲}

\postdate{2022年12月3日}

\result{\frac{1}{2}\qty(x^2\arctan{x}-x+\arctan{x})+C}

\analysis{
    分析:在我们现有的三大方法（凑微分法，换元法和分部积分法）中，这题凑微分法和换元法是很难操作的，因为感觉不到其中凑的关系.而对于这样的内容，我们考虑分部积分法.


    上一期的每日一题里我们介绍了升阶降阶的思想，也就是想办法把比较麻烦，不好积 分的东西转化为求导.那么今天我们介绍一个经验公式：反对幂指三.


    先简单介绍一下这句话的意思：反三角函数<对数函数<幂函数<指数函数<三角函数


    这个经验公式其实是根据这些函数求积分的难度和原函数形式的复杂程度由高到低排列的.那么我们应该怎么使用呢？


    其实就是在使用不定积分的时候，我们可以把“反对幂指三”中靠右的部分优先升阶“收”到$\d$后面，会方便计算.比如我们今天的这一题.


    根据这个经验公式，$x$更靠右，我们应该将$x$放到$\d$后面去.
    \formula{
        \begin{aligned}
            \int{x\arctan x\d x}&=\frac{1}{2}\int{\arctan{x}\d x^2}\\
            &=\frac{1}{2}\qty(x^2\arctan{x}-\int{x^2\d\arctan{x}})\\
            &=\frac{1}{2}\qty(x^2\arctan{x}-\int{\frac{x^2}{x^2+1}\d x})\\
            &=\frac{1}{2}\qty(x^2\arctan{x}-\int{1-\frac{1}{x^2+1}\d x})\\
            &=\frac{1}{2}\qty(x^2\arctan{x}-x+\arctan{x})+C
        \end{aligned}
    }
    不过这个终究是经验公式，而且在实际操作中，反三角函数$\approx$对数函数，指数函数$\approx$三角函数，因此也会有“反对幂三指”的说法.同时，这只是经验公式，是大部分情况是正确的，但不代表绝对的正确.具体情况还是要具体分析.


    比如$\int{\tan{x}\d x}=x\tan{x}-\int{x\d\tan{x}}$


    操作起来并不如$\int{\tan{x}\d x}=\int{\frac{\sin{x}}{\cos{x}}\d x}=-\int{\frac{\d\cos{x}}{\cos{x}}}$方便.
}